\documentclass[12pt]{article}
\usepackage{amssymb}
\usepackage{template}

\homeworknumber{1}

\begin{document}

\makehomeworkheader
\problem{PCA for Different Distributions}
\begin{subquestions}
\item Write out the PMF of a negative binomial distribution. When would you use a negative binomial distribution versus a Poisson distribution?

% \vspace{3cm}

\item Write out the GLM-PCA optimization problem for single cell count data following a negative binomial distribution. (Hint: The negative binomial distribution has two arguments, mean and overdispersion. Only use the matrix factorization $UV$ for one of these two arguments, and set the other as a learned constant.)

% \vspace{3cm}

\item What is the PDF of a log-normal distribution? When would you use such a distribution?

% \vspace{3cm}

\item Suppose your collaborator generated metabolomics data and wanted you to run GLM-PCA on it. After preprocessing, metabolomics data is often approximated with a log-normal distribution. Write out the GLM-PCA optimization problem for this metabolomics data.

% \vspace{3cm}

\item Suppose you asked ChatGPT to answer Question 1d for a brand new, very complicated distribution, and it gave you a specific optimization problem. However, you did not trust ChatGPT's answer and wanted to double-check that it gave you the correct optimization problem. Suggest a way you could do so without re-deriving the optimization problem from scratch.

% \vspace{3cm}
\end{subquestions}
\newpage

\problem{Canonical correlation analysis (CCA)}

We are given two single cell datasets
\begin{equation}
    X=[x_{ig}] \in \mathbb{R}^{N \times G} \qquad \text{and} \qquad Y=[y_{ip}] \in \mathbb{R}^{N \times P},
\end{equation}
where each dataset measures a molecular feature (e.g., gene expression, protein expression, chromatin accessibility) on $N$ cells. The two datasets measure different features, with sizes $G$ and $P$ respectively, on the same $N$ cells. Assume $X$ and $Y$ are mean-centered, meaning that $\frac{1}{N} \sum_{i=1}^N x_{ig} = 0$ and $\frac{1}{N} \sum_{i=1}^N y_{ip} = 0$ for all features $g=1,\dots,G$ and $p=1,\dots,P$.

Canonical correlation analysis (CCA) is an approach to identify ``shared correlation structures'' for the two datasets $X$ and $Y$. Specifically, in CCA, one aims to identify linear combinations
\begin{equation}
u \in \mathbb{R}^{G} \qquad \text{and} \qquad v \in \mathbb{R}^{P},    
\end{equation}

of the two features such that $Xu$ and $Yv$ are maximally correlated.

\begin{subquestions}
\item Fill in the blanks in the CCA optimization problem:
\begin{equation}
    \argmax_{\substack{||u||_2=1 \\ ||v||_2=1}} \;\Corr\left(\underline{\hspace{1.5cm}},\underline{\hspace{1.5cm}} \right)
\end{equation}


% \underset{\lVert u\rVert_2^2=1,\ \lVert v\rVert_2^2=1}{\operatorname{argmax}}
% \operatorname{Corr}\left(\underline{\hspace{1.5cm}},\underline{\hspace{1.5cm}}\right).
% \]

\item Suppose $X=Y$ and we replace $\operatorname{Corr}$ with $\operatorname{Cov}$ in the optimization problem in part a. What problem does this represent? (Hint: We spent a class on it.)

\item Recall the formula for Pearson correlation between two vectors $a$ and $b$:
\begin{equation}
    \Corr(a,b) = \frac{\Cov(a,b)}{\sqrt{\Var(a) \Var(b)}}
\end{equation}
% \[
% \operatorname{Corr}(a,b)
% =
% \frac{\operatorname{Cov}(a,b)}
% {\sqrt{\operatorname{Var}(a)\operatorname{Var}(b)}}.
% \]

Rewrite the optimization problem in part \textbf{a} using this equation.

\item Define the variables $\Sigma_{XY}=\frac{1}{N-1}X^{T}Y$, $\Sigma_X=\frac{1}{N-1}X^{T}X$, and $\Sigma_Y=\frac{1}{N-1}Y^{T}Y$.

Rewrite the optimization problem in part \textbf{c} to remove the variables $X$ and $Y$, and instead use the variables $\Sigma_{XY}, \Sigma_X$, and $\Sigma_Y$.

The variable $\Sigma_{XY}$ is called the \emph{sample cross-covariance matrix} and the variables $\Sigma_X$ and $\Sigma_Y$ are called the \emph{sample covariance matrices} for $X$ and $Y$, respectively.

Hint: Your optimization problem should have the form
\begin{equation}
    \argmax_{\substack{||v||_2=1, ||w||_2=1}} \frac{}{\sqrt{\hspace{2.5cm}}
 \sqrt{\hspace{2.5cm}}}
\end{equation}
% \[
% \underset{\lVert u\rVert_2=1,\ \lVert v\rVert_2=1}{\operatorname{argmax}}
% \frac{\hspace{2.5cm}}
% {\sqrt{\hspace{2.5cm}}
%  \sqrt{\hspace{2.5cm}}}.
% \]

\item Define the variables 
\begin{equation}
    \widetilde{u}=\Sigma_X^{1/2}u
\qquad\text{and}\qquad
\widetilde{v}=\Sigma_Y^{1/2}v.
\end{equation}

Rewrite the optimization problem in part \textbf{d} in terms of the variables $\widetilde{u}, \widetilde{v}, \Sigma_X, \Sigma_Y,\Sigma_{XY}$.

\item Suppose you are told that a solution to the optimization problem
\begin{equation}
    \argmax_{\substack{||a||_2=1 \\ ||b||_2=1}} a^T M b
\end{equation}
is given by
\begin{equation}
    \widehat{a}=u
\qquad\text{and}\qquad
\widehat{b}=v,
\end{equation}
where $u$ and $v$ are the (matched) leading left and right singular vectors of the matrix $M$. 

Use this knowledge to derive the solution to the optimization problem in part \textbf{e}?

\item \textit{[Open-ended]} The procedure you derived gives you the first \emph{canonical correlation vectors} $u_1$ and $v_1$. How could you extend the approach in part \textbf{f} to find the second canonical correlation vectors $u_2$ and $v_2$? More generally, how could you find the $k$-th canonical correlation vectors $u_k$ and $v_k$ for any $k$?

\item Consider a different setting where we are given two single cell datasets
\begin{equation}
    X \in \mathbb{R}^{N \times G}
\qquad \text{and} \qquad
Y \in \mathbb{R}^{M \times G}
\end{equation}
which measure the same molecular feature (e.g., gene expression) on two cell populations of sizes $N$ and $M$, respectively.

We would like to use CCA to embed the $N$ cells and $M$ cells, respectively, into a shared embedding space. How could you use CCA to do this?
\end{subquestions}

\textit{[CCA is the core algorithmic idea behind the first version of \href{https://www.nature.com/articles/nbt.4096}{Seurat}, a very popular single cell analysis package.]}

\newpage

\problem{NMF for multiple samples and modalities}

% Suppose we collect non-negative single cell data from $S$ biological samples using $M$ different assays/modalities. For each sample $s=1,\dots,S$ and modality $m=1,\dots,M$, let $X^{(s,m)} \in \mathbb{R}^{N_s \times P_m}$ be the corresponding data matrix of size $N_s \times P_m$, where $N_s$ is the number of cells in sample $s$ and $P_m$ is the number of features measured by modality $m$.

% For simplicity, we assume that w runs several different modalities on the same sample, they are running it on the same $N_s$ cells.

% We aim to factorize each observed data matrix $X^{(s,m)}$ as $X^{(s,m)} \approx U^{(s)}V^{(m)}$, where $U^{(s)} \in \mathbb{R}^{N_s \times K}$ is the latent factor representations of the cells in sample $s$ and $V^{(m)} \in \mathbb{R}^{K \times P_m}$ is the feature loadings for a modality $m$.

% \begin{enumerate}[label=\textbf{(\alph*)}, leftmargin=*]

\begin{subquestions}
\item
Suppose we measure a single cell data matrix $X \in \mathbb{R}^{N \times P}$ for $N$ cells with $P$ features. We aim to factorize the observed data as $X \approx UV$, where $U \in \mathbb{R}^{N\times K}$ is the factor matrix and $V\in \mathbb{K \times G}$ is the loading matrix, for some small $K>0$.
% one sample and one modality, i.e. $S=M=1$, so that the data consists of only a single matrix $X \in \mathbb{R}^{N \times P}$.
Write out the NMF optimization problem for computing this factorization.
% Qualitatively explain what the rows of $U$ and $V$ represent.


\item 
% \textbf{Multiple samples, one modality}
Next, suppose we measure $S$ different single cell data matrices $X^{(s)} \in \mathbb{R}^{N_s \times P}$ for $s=1,\dots,S$. Our data measures the same $P$ features across $S$ different \emph{samples}, i.e. populations of cells.
% $S>1$ samples for a single modality (i.e. $M=1$). Then, the data consists of matrices $X^{(s)} \in \mathbb{R}^{N_s \times P}$ which measure the same $P$ features across $S$ different populations of cells. 
We aim to factorize each sample as $X^{(s)} \approx U^{(s)} V$, where each sample $s$ has its own factor matrix $U^{(s)}$ but all samples share a loading matrix $V$.

% Assume that each sample will have its own cell-factor matrix $U^{(s)}$, but all samples will share one gene-loading matrix $V$: $X^{(s)} \approx U^{(s)}V$.

Write out an NMF-style optimization problem for computing this factorization. Qualitatively explain what biological assumption is made if we share the loading matrix $V$ across all samples.

\item
% \item \textbf{Multiple modalities, same sample}
Next, suppose we measure $M$ different single cell data matrices $X^{(m)} \in \mathbb{R}^{N \times P_m}$ for $m=1,\dots,M$. Our data measures $M$ different molecular modalities (e.g. gene expression, chromatin accessibility, etc.) across the same $N$ cells.
% $S>1$ samples for a single modality (i.e. $M=1$). Then, the data consists of matrices $X^{(s)} \in \mathbb{R}^{N_s \times P}$ which measure the same $P$ features across $S$ different populations of cells. 
We aim to factorize each sample as $X^{(m)} \approx U V^{(m)}$, where each modality $m$ has its own loading matrix $V^{(m)}$ but all modalities share a factor matrix $U$.

Write out an NMF-style optimization problem for computing this factorization. Qualitatively explain what biological assumption is made if we share the factor matrix $U$ across all samples.


% Now, suppose a single set of $N$ cells is measured using $M$ modalities: $X^{(m)} \in \mathbb{R}_{+}^{N \times G_m}$, $m = 1, \ldots, M$. Here the modalities share the same cell-factor matrix $U$, but each modality has its own feature loading matrix $V^{(m)}$: $X^{(m)} \approx UV^{(m)}$.

% Write out the joint NMF optimization problem. Is it reasonable to share $U$ and/or $V^{(m)}$ across modalities? Why/why not? If the modalities have different scales for the features, what do we have to take into consideration?



% \item \textbf{Multiple modalities, multiple samples}

\item
Next, suppose we measure $S\cdot M$ different single cell data matrices $X^{(s,m)} \in \mathbb{R}^{N_s \times P_m}$ for $s=1,\dots,S$ and $m=1,\dots,M$. Our data measures $M$ different molecular modalities across $S$ different samples.
We aim to factorize each sample as $X^{(s,m)} \approx U^{(s)} V^{(m)}$, where each sample $s$ has its own factor matrix $U^{(s)}$ and each modality $m$ has its own loading matrix $V^{(m)}$.

Write out an NMF-style optimization problem for computing this factorization. Qualitatively explain what biological assumption(s) are being made by this factorization.


% Now suppose all $M$ modalities are measured in all $S$ samples. Each sample has its own cell-factor matrix $U^{(s)}$ and each modality has its own feature-loading matrix $V^{(m)}$: $X^{(s,m)} \approx U^{(s)}V^{(m)}$.

% Write a single objective function to jointly model all $S \times M$ data matrices. By looking at the function, can you tell which information is shared across samples, and which information is shared across modalities?



% \item \textbf{Missing some sample-modalities pairs}
\item
In practice, not every modality may be measured in every sample.
Let $\Omega \subseteq \{1, \ldots, S\} \times \{1, \ldots, M\}$ be the set of observed pairs $(s,m)$ of samples and modalities. sample-modality pairs. We only observe single cell data matrices $X^{(s,m)}$ for sample-modality pairs $(s,m) \in \Omega$. 

Modify the optimization problem in part \textbf{d} to only incorporate the observed data.

% Because $X^{(s,m)}$ is unknown when $(s,m) \notin \Omega$, we can’t calculate its reconstruction error. Instead, modify the objective function from part (d) such that only observed pairs contribute to the loss.



\item
Assume that observed pairs are sufficient to learn every factor matrix $U^{(s)}$ and loading matrix $V^{(m)}$. How would you estimate an \emph{unobserved} $X^{(s,m)}$ (i.e. $(s,m) \notin \Omega$)?



% Of course, in practice not every modality may have been measured in every sample. We see a version of this in \href{https://doi.org/10.1186/s13059-020-02015-1}{MOFA+ (Argelaguet et al., 2020)}, a widely used probabilistic factor model for integrating multi-omic data that can accommodate missing modalities. However in this question we will simplify the setting and use the NMF model above, assuming there is only one project containing $S$ samples and $M$ possible modalities.

% Now, let $\Omega \subseteq \{1, \ldots, S\} \times \{1, \ldots, M\}$ represent the observed existing sample-modality pairs. Because $X^{(s,m)}$ is unknown when $(s,m) \notin \Omega$, we can’t calculate its reconstruction error. Instead, modify the objective function from part (d) such that only observed pairs contribute to the loss.
% \vspace{2.5cm}

% Now, assume that observed pairs are sufficient to learn every $U^{(s)}$ and $V^{(m)}$. How would you estimate a missing matrix $X^{(s,m)}$ for which $(s,m) \notin \Omega$? (Hint: use the learned $U^{(s)}$ and $V^{(m)}$ matrices!)



% \end{enumerate}

\end{subquestions}



% \end{enumerate}

\end{document}
